% ============================================================================
% QFA-LTL: Verification of Noisy Quantum Executions
% IEEE QCE 2026 submission skeleton
% ============================================================================
\documentclass[conference]{IEEEtran}
\usepackage{amsmath,amssymb,amsfonts}
\usepackage{graphicx}
\usepackage{booktabs}
\usepackage{multirow}
\usepackage{algorithmic}
\usepackage{url}
\usepackage{cite}

\title{Verification of Noisy Quantum Executions:\\ A Derived, Calibration-Free
Adaptive Anchor with a Characterized Detectability Frontier}

\author{\IEEEauthorblockN{Y. Agnihotri}
\IEEEauthorblockA{Research project, IEEE QCE 2026 submission}}

\begin{document}
\maketitle

% ---------------------------------------------------------------------------
\begin{abstract}
Verifying a quantum circuit on NISQ hardware raises a question that classical
equivalence checking does not answer: \emph{is this run's measured behavior
acceptable against its temporal specification, given the device's current
noise?} Static thresholds fail on bad calibration days (false failures) and
miss faults on good days. We present an oracle-synthesis framework whose
pass/fail threshold is derived from first principles:
\begin{equation}
\tau = p_{\mathrm{ideal}} \cdot (1-\varepsilon)^V - \mathrm{Wilson}_{95\%}(n_{\mathrm{shots}}),
\label{eq:anchor}
\end{equation}
where $p_{\mathrm{ideal}}$ is the exact noise-free probability of the target
basis state(s) from a validated statevector compiler
(max $|\Delta p| = 8.9\times10^{-16}$), $\varepsilon$ is the loss measured by
a Bell-state calibration probe on the circuit's mapped physical qubits, and
$V$ is the two-qubit-gate volume. The anchor is calibration-free by
construction: no outcome data from the circuits under test is used.

Over 10 seeds with a 60/40 hold-out split on a 24-circuit suite under
Brisbane calibration noise, the framework achieves \textbf{0.0\% false-failure
rate}, \textbf{100\% alarm precision}, and 67\% bug-detection recall. The
undetected faults are provably measurement-invisible (they preserve the
target-basis probability distribution). A calibration-drift simulation shows
the anchor tracking device noise where a fixed $0.15$ threshold false-fails
correct circuits. A comparison against QCEC (17/18 ideal-circuit detection)
documents the complementary roles of ideal-unitary correctness versus
noisy-run acceptability.
\end{abstract}

% ---------------------------------------------------------------------------
\section{Introduction}
\label{sec:intro}
% Problem: NISQ verification is about noisy measured behavior, not ideal
% unitaries. Equivalence checking answers "is the circuit right?"; this work
% answers "is this run acceptable?" — a question equivalence checking cannot
% answer, because it operates on noise-free semantics.
% Contributions: (i) derived anchor \eqref{eq:anchor}, (ii) validated exact
% compiler with truncation bound, (iii) detectability frontier
% characterization, (iv) open benchmark suite + fault taxonomy.
% Related work: quantum program testing (metamorphic testing, oracle
% problem), equivalence checking (QCEC/MQT), statistical verification,
% LTL monitoring (classical), QFA naming caveat (our automaton is a
% classical weighted transition system — disambiguate explicitly).

\section{Preliminaries and Related Work}
\label{sec:related}
% LTL and monitor automata; windowed Bernoulli predicates.
% Quantum noise models: depolarizing, readout, calibration data.
% QCEC / decision diagrams: ideal-unitary correctness.
% Statistical verification: concentration bounds, Wilson intervals.
% Position: the adaptive anchor is a *test oracle* for noisy runs.

\section{Framework}
\label{sec:framework}
\subsection{Specification language}
% Grammar: F/G/F<=k/G<=k ( prob(|b1>)+...+prob(|bm>) c t ).
% Windowed Bernoulli semantics; F vs G acceptance; violating states.

\subsection{The derived anchor}
\label{sec:anchor}
% Derivation of \eqref{eq:anchor} from a per-gate depolarizing model:
% target probability decays as (1-eps)^V; observed rate is a binomial
% sample of q = p_ideal (1-eps)^V; Wilson lower bound gives the threshold.
% eps measured by the ZKC Bell probe (4096 shots on mapped qubits).
% Calibration-free: p_ideal from the compiler, eps from the device.

\subsection{Exact circuit-to-automaton compiler}
% Statevector tracking; gate handlers; validated vs Qiskit (8.9e-16);
% truncation with tracked error bound.

\subsection{Verdict pipeline}
% Metrics: probability / parity / expectation. Direction from spec.
% Optional gray zone: borderline = needs review.
% Seeded execution; structured reports.

\section{Benchmark Suite and Protocol}
\label{sec:benchmark}
% 6 algorithms x (correct + 3 faults) = 24 circuits; fault catalog.
% Brisbane noise model via Aer; 1024 shots; 3 reps; 10 seeds;
% 60/40 hold-out split.
% Metrics (buggy = positive): recall, precision, false-failure rate.

\section{Results}
\label{sec:results}
\begin{table}[t]
\centering
\caption{Strategy comparison (10 seeds, hold-out 60/40; mean over seeds).}
\label{tab:strategies}
\begin{tabular}{lccc}
\toprule
Strategy & FF rate & Recall & Precision \\
\midrule
derived (ours) & \textbf{0.000} & 0.670 & \textbf{1.000} \\
fixed 0.15     & 0.167 & 0.550 & 0.908 \\
Hoeffding      & 0.000 & 0.428 & 1.000 \\
Wilson         & 0.000 & 0.528 & 1.000 \\
no-ZKC ablation & 0.967 & 1.000 & 0.757 \\
\bottomrule
\end{tabular}
\end{table}
% R1: derived anchor: 0% FF, 100% precision, F1 0.799 (highest among
%     principled strategies).
% R2: McNemar: derived vs no_zkc p=0.0009 (probe decisive); vs fixed ns
%     (tradeoff axis: zero FF at slightly lower recall).
% R3: sensitivity sweep: wide plateau of perfect recall; precision 1.0.
% R4: compiler validation: max |dp| = 8.9e-16.
% R5: hold-out robustness: calibration-free design does not degrade.
% R6: drift study: 2/8 days recovered where fixed false-fails.

\section{The Detectability Frontier}
\label{sec:frontier}
% Proposition sketch: if fault U preserves |<b|U|0>|^2 for all target-
% relevant b, no measurement-statistics oracle can detect it. Empirical
% confirmation: QFT phase/swap (uniform spectrum preserved), GHZ
% entanglement mutations (parity preserved), QAOA wrong-problem.
% QCEC comparison: 17/18 detected on ideal circuits — documents the
% complementarity, and validates the suite (faults are real).
% Implication: phase-sensitive probes (rotated bases, entanglement
% witnesses) are the future work.

\section{Discussion and Limitations}
\label{sec:discussion}
% Simulator-validated; real-device drift study is the decisive next step.
% The anchor's statistical guarantee structure (Wilson margin) vs a full
% soundness theorem (Phase 2).
% Gray-zone tradeoff: protects correct circuits at recall cost.
% Scope: not equivalence checking, not post-classical verification.

\section{Conclusion}
\label{sec:conclusion}

% ---------------------------------------------------------------------------
\section*{Acknowledgment}
% None.

\bibliographystyle{IEEEtran}
% \bibliography{refs}

\end{document}
